% bismi ALLAH ARRAHMAN ARRAHIM






























































\documentclass[11pt,a4paper]{article}

% -------------------------------------------------
% Encodage, langue, typographie
% -------------------------------------------------
\usepackage[french]{babel}
\usepackage{fontspec}
\usepackage{microtype}
\usepackage[a4paper,margin=2.5cm]{geometry}

% -------------------------------------------------
% Mathématiques
% -------------------------------------------------
\usepackage{amsmath,amssymb,amsthm,mathtools,bm}
\usepackage{enumitem}

% -------------------------------------------------
% Liens
% -------------------------------------------------
\usepackage[hidelinks]{hyperref}

% -------------------------------------------------
% Police
% -------------------------------------------------
\setmainfont{TeX Gyre Pagella}

% -------------------------------------------------
% Environnements
% -------------------------------------------------
\newtheorem{definition}{Définition}[section]
\newtheorem{assumption}{Hypothèse}[section]
\newtheorem{proposition}{Proposition}[section]
\newtheorem{theorem}{Théorème}[section]
\newtheorem{corollary}{Corollaire}[section]
\newtheorem{remark}{Remarque}[section]

% -------------------------------------------------
% Opérateurs
% -------------------------------------------------
\DeclareMathOperator{\prox}{prox}
\DeclareMathOperator{\argmin}{argmin}
\DeclareMathOperator{\supp}{supp}
\DeclareMathOperator{\KL}{KL}
\DeclareMathOperator{\Id}{Id}
\DeclareMathOperator{\diag}{diag}
\DeclareMathOperator{\softmax}{softmax}

% -------------------------------------------------
% Raccourcis
% -------------------------------------------------
\newcommand{\X}{\mathcal X}
\newcommand{\C}{\mathcal C}
\newcommand{\D}{\mathcal D}
\newcommand{\E}{\mathcal E}
\newcommand{\Z}{\mathcal Z}
\newcommand{\Hh}{\mathcal H}
\newcommand{\I}{\mathcal I}
\newcommand{\T}{\mathcal T}
\newcommand{\W}{\mathcal W}
\newcommand{\N}{\mathcal N}
\newcommand{\R}{\mathbb R}
\newcommand{\Pcal}{\mathcal P}
\newcommand{\Qcal}{\mathcal Q}
\newcommand{\Gcal}{\mathcal G}
\newcommand{\Mcal}{\mathcal M}
\newcommand{\1}{\mathbf 1}

\title{PCH-ESM\\
\large Purification Contextuelle Hiérarchique et Espaces Sémantiques Modulaires}
\author{}
\date{}

\begin{document}

\maketitle

\begin{abstract}
Ce document formalise une architecture hiérarchique pour modèles séquentiels et systèmes de langage, fondée sur trois piliers :
(i) la purification récursive d'une représentation dans un contexte donné,
(ii) la décomposition hiérarchique d'un espace sémantique global en sous-espaces spécialisés,
et (iii) une fonction de coordination \(G\) qui transfère la structure de trajectoire d'un modèle généraliste vers un ensemble de modèles spécialisés.
La théorie distingue explicitement la représentation, la stabilité contextuelle, la hiérarchie, la spécialisation, la coordination inter-espaces et le rôle de modèle du monde.
\end{abstract}

\tableofcontents

% =================================================
\section{Cadre général}
% =================================================

L'idée directrice de PCH-ESM est qu'une bonne représentation n'est pas seulement une représentation prédictive, mais une représentation stable \emph{dans un contexte}.
Cette stabilité n'est pas absolue : elle dépend de la tâche, du domaine, du mode d'exécution, des contraintes de calcul et de la structure locale de coopération entre sous-modèles.

On distingue donc trois niveaux :

\begin{enumerate}[leftmargin=2em]
    \item un niveau de purification hiérarchique de base, noté PH ;
    \item un niveau contextuel, noté PCH ;
    \item un niveau architectural distribué, noté Multi-ES.
\end{enumerate}

L'architecture complète combine ensuite une couche de coordination \(G\), puis éventuellement une couche de modèle du monde Multi-PH.

% =================================================
\section{PH : purification hiérarchique}
% =================================================

Cette section résume la base PH déjà présente dans le document EEG.

\begin{definition}[Opérateur de purification hiérarchique]
Soit \(u\) un objet séquentiel ou latent. On définit un opérateur de purification \(\Pcal_\alpha\), de paramètre \(\alpha\in[0,1]\), qui agit récursivement via deux passes directionnelles :
\[
\mathcal F_\alpha(u)_i
=
\alpha\,\mathcal F_\alpha(u)_{i-1}
+
(1-\alpha)u_i,
\]
\[
\mathcal B_\alpha(u)_i
=
\alpha\,\mathcal B_\alpha(u)_{i+1}
+
(1-\alpha)u_i,
\]
et la purification symétrique :
\[
\Pcal_\alpha(u)
=
\frac12\Bigl(\mathcal F_\alpha(u)+\mathcal B_\alpha(u)\Bigr).
\]
\end{definition}

\begin{definition}[Hiérarchie de noyaux]
On définit :
\[
N^{(0)} = W,
\qquad
N^{(r+1)} = \Pcal_\alpha\!\left(N^{(r)}\right),
\quad r\ge 0.
\]
Le noyau terminal est
\[
N^{(\star)} = N^{(r^\star)}
\]
lorsque la suite converge ou lorsqu'une profondeur maximale est atteinte.
\end{definition}

\begin{definition}[Résidu]
Le résidu associé au noyau terminal est
\[
R^{(0)} = W - N^{(\star)}.
\]
Il peut être purifié à son tour, possiblement sous guidage du noyau terminal.
\end{definition}

% =================================================
\section{De PH vers PCH : purification contextuelle hiérarchique}
% =================================================

Cette section généralise PH en introduisant explicitement un contexte.

\subsection{Espace des contextes}

\begin{definition}[Espace des contextes]
On note \(\C\) l'espace des contextes.
Un contexte est un élément
\[
c\in\C.
\]
On peut décomposer un contexte sous la forme
\[
c=(d,\tau,m,\omega,\xi),
\]
où :
\begin{itemize}[leftmargin=2em]
    \item \(d\) désigne le domaine ou la spécialité ;
    \item \(\tau\) désigne la tâche ou l'objectif ;
    \item \(m\) désigne le mode d'exécution ou de coopération ;
    \item \(\omega\) désigne des contraintes de calcul, mémoire ou temps ;
    \item \(\xi\) désigne des paramètres additionnels propres à l'application.
\end{itemize}
\end{definition}

\begin{remark}
Le contexte n'est pas une simple étiquette. Il encode la manière dont l'objet doit être interprété, purifié et utilisé.
Deux contextes distincts peuvent donc induire deux stabilités différentes pour une même représentation.
\end{remark}

\subsection{Stabilité contextuelle}

\begin{definition}[Fonctionnelle de stabilité contextuelle]
Soit \(\Z\) un espace latent.
On définit une fonctionnelle
\[
J:\Z\times\C\to\R\cup\{+\infty\},
\]
où \(J(z;c)\) mesure le coût d'instabilité de la représentation \(z\) dans le contexte \(c\).
On suppose que, pour chaque \(c\in\C\), la fonction \(z\mapsto J(z;c)\) est propre, semi-continue inférieurement, et au moins coercive sur le sous-espace considéré.
\end{definition}

\begin{definition}[Stabilité contextuelle]
Une représentation \(z^\star\in\Z\) est dite \emph{stable dans le contexte \(c\)} si
\[
z^\star\in\argmin_{z\in\Z} J(z;c).
\]
Si l'ensemble des minimisateurs est singleton, on parle de \emph{noyau contextuel unique}.
Sinon, on parle d'un \emph{ensemble de noyaux contextuels}.
\end{definition}

\begin{remark}
La stabilité n'est donc plus globale ni absolue : elle est relative au contexte \(c\).
C'est précisément ce qui distingue PCH de PH.
\end{remark}

\subsection{Opérateur de purification contextuelle}

\begin{definition}[Opérateur de purification contextuelle]
Pour un pas \(\lambda>0\), on définit
\[
\Pcal_{\lambda}(\,\cdot\,;c):\Z\to\Z,
\qquad
\Pcal_{\lambda}(z;c)
=
\argmin_{u\in\Z}
\left\{
J(u;c)
+
\frac{1}{2\lambda}\|u-z\|^2
\right\}.
\]
\end{definition}

\begin{remark}
Lorsque \(J(\cdot;c)\) est convexe, \(\Pcal_{\lambda}(\cdot;c)\) est l'opérateur proximal associé à la stabilité contextuelle.
Lorsque \(J\) est non convexe, la définition reste formelle mais nécessite un choix de branche ou une approximation numérique.
\end{remark}

\subsection{Trajectoires de purification contextuelle}

\begin{definition}[Trajectoire de purification contextuelle]
Soit une entrée \(x\in\X\), un encodeur \(E_\theta:\X\to\Z\), et un contexte \(c\in\C\).
La trajectoire de purification contextuelle associée à \((x,c)\) est la suite
\[
z^{(0)}(x,c)=E_\theta(x),
\qquad
z^{(r+1)}(x,c)=\Pcal_{\lambda}\!\bigl(z^{(r)}(x,c);c\bigr),
\quad r\ge 0.
\]
On note cette trajectoire
\[
\T(x,c)=\bigl(z^{(0)}(x,c),z^{(1)}(x,c),z^{(2)}(x,c),\dots\bigr).
\]
\end{definition}

\begin{definition}[Noyau contextuel]
Si la suite \(\{z^{(r)}(x,c)\}_{r\ge 0}\) converge, on définit le noyau contextuel par
\[
N^\star(x,c)=\lim_{r\to\infty}z^{(r)}(x,c).
\]
Si la convergence numérique n'est pas garantie, on prend une approximation au rang \(r^\star\) :
\[
N^\star(x,c)\approx z^{(r^\star)}(x,c).
\]
\end{definition}

\begin{definition}[Résidu contextuel]
Le résidu contextuel associé à \((x,c)\) est
\[
R^{(0)}(x,c)=z^{(0)}(x,c)-N^\star(x,c).
\]
Il représente la partie de la représentation initiale qui n'est pas absorbée par la stabilité contextuelle principale.
\end{definition}

\subsection{PH comme cas limite de PCH}

\begin{theorem}[PH est un cas limite de PCH]
Supposons que le contexte soit trivial, c'est-à-dire qu'il existe un unique contexte \(c_0\in\C\), ou bien que la fonctionnelle de stabilité ne dépende pas effectivement du contexte :
\[
J(z;c)=J(z)
\quad\text{pour tout } c\in\C.
\]
Alors l'opérateur de purification contextuelle devient indépendant du contexte et se réduit à l'opérateur de PH :
\[
\Pcal_{\lambda}(z;c)=\Pcal_{\lambda}(z),
\]
et la trajectoire contextuelle se confond avec la trajectoire PH classique.
\end{theorem}

\begin{proof}
Si \(\C=\{c_0\}\), tout objet contextuel devient constant en \(c\), donc la dépendance contextuelle disparaît par trivialité.
Si, plus généralement, \(J(\cdot;c)\) est indépendant de \(c\), alors pour tout \(z\in\Z\),
\[
\Pcal_{\lambda}(z;c)
=
\argmin_{u\in\Z}
\left\{
J(u)+\frac{1}{2\lambda}\|u-z\|^2
\right\},
\]
ce qui ne dépend plus de \(c\).
On retrouve donc exactement l'opérateur de purification de PH.
\end{proof}

\begin{corollary}
PH s'obtient comme le sous-cas de PCH dans lequel le contexte est dégénéré, constant ou non discriminant.
Ainsi, PCH généralise PH.
\end{corollary}

\subsection{Version bidirectionnelle contextuelle}

Pour rester fidèle à la purification bidirectionnelle de PH, on peut introduire une version contextuelle symétrique.

\begin{definition}[Purification bidirectionnelle contextuelle]
Soit \(u=(u_1,\dots,u_T)\) une séquence latente locale.
On définit :
\[
\mathcal F_{\alpha,c}(u)_i
=
\alpha(c)\,\mathcal F_{\alpha,c}(u)_{i-1}
+
\bigl(1-\alpha(c)\bigr)u_i,
\]
\[
\mathcal B_{\alpha,c}(u)_i
=
\alpha(c)\,\mathcal B_{\alpha,c}(u)_{i+1}
+
\bigl(1-\alpha(c)\bigr)u_i,
\]
où \(\alpha:\C\to[0,1]\) est une fonction de contrôle contextuel.
La purification contextuelle bidirectionnelle est alors
\[
\Pcal_{\alpha,c}(u)
=
\frac12\Bigl(\mathcal F_{\alpha,c}(u)+\mathcal B_{\alpha,c}(u)\Bigr).
\]
\end{definition}

\begin{remark}
Dans cette version, le contexte influence le degré de lissage, la profondeur effective de récursion, et le compromis entre compression et reconstruction.
\end{remark}

% =================================================
\section{PCH-ESM : espaces sémantiques modulaires et hiérarchiques}
% =================================================

Cette section formalise la décentralisation de l'espace sémantique, la spécialisation locale, la coordination globale, la gestion de la mémoire et du calcul, ainsi que l'émergence de trajectoires de pensée traçables.

\subsection{Définition générale}

\begin{definition}[Système PCH-ESM]
Un système PCH-ESM est un tuple
\[
\mathfrak M
=
\Bigl(
\E_G,
\{\E_i\}_{i=1}^n,
\{\Pi_i,\iota_i,\Pcal_i,\pi_i\}_{i=1}^n,
\Gcal,
r,
\Gamma
\Bigr),
\]
où :

\begin{itemize}[leftmargin=2em]
    \item \(\E_G\) est l'espace sémantique global ;
    \item \(\E_i\) est le \(i\)-ème espace sémantique spécialisé ;
    \item \(\Pi_i:\E_G\to\E_i\) est la projection locale ;
    \item \(\iota_i:\E_i\to\E_G\) est l'injection ou remontée vers l'espace global ;
    \item \(\Pcal_i\) est l'opérateur local de purification dans \(\E_i\) ;
    \item \(\pi_i\) est l'agent ou politique locale responsable de \(\E_i\) ;
    \item \(\Gcal\) est un graphe orienté acyclique de dépendances entre espaces ;
    \item \(r\) est une politique de routage ;
    \item \(\Gamma\) est l'opérateur de fusion globale.
\end{itemize}
\end{definition}

\subsection{Compatibilité entre espace global et espaces locaux}

\begin{assumption}[Compatibilité des projections]
Pour tout \(i\in\{1,\dots,n\}\),
\[
\Pi_i\circ \iota_i = \Id_{\E_i}.
\]
De plus, pour tout \(z\in\E_G\), l'assemblage des composantes locales reconstruit approximativement l'état global :
\[
\left\| z - \Gamma\bigl(\iota_1\Pi_1 z,\dots,\iota_n\Pi_n z\bigr)\right\|
\le \varepsilon_{\mathrm{coh}}.
\]
\end{assumption}

\begin{remark}
Cette condition exprime le fait qu'un espace global peut être décomposé en sous-espaces spécialisés sans perdre toute cohérence, à condition que la fusion et les projections soient compatibles.
\end{remark}

\subsection{Graphe hiérarchique de collaboration}

On considère un graphe orienté acyclique
\[
\Gcal=(\I,\mathcal E),
\]
dont les sommets sont les espaces spécialisés et dont les arêtes représentent des canaux de communication, de dépendance ou de transfert de contexte.

\begin{definition}[Trajectoire de pensée]
Une trajectoire de pensée est une suite finie
\[
\tau
=
\bigl((i_1,z_1,m_1),(i_2,z_2,m_2),\dots,(i_T,z_T,m_T)\bigr),
\]
où \(i_t\in\I\) est l'espace actif à l'instant \(t\), \(z_t\in\E_{i_t}\) l'état local, et \(m_t\) un message ou une mise à jour émise vers d'autres espaces.
\end{definition}

\begin{remark}
Cette définition rend la trajectoire de calcul explicitement observable. Le modèle n'est plus seulement une boîte noire qui produit une sortie ; il devient un système dont le cheminement interne peut être inspecté.
\end{remark}

\subsection{Routage et activation sparse}

\begin{definition}[Politique de routage]
La politique de routage est une application
\[
r:\E_G\times\C\to\Delta(\I),
\]
où \(\Delta(\I)\) désigne le simplexe des distributions sur \(\I\).
Pour un état global \(z\) et un contexte \(c\), la politique \(r(z,c)\) donne la probabilité ou la pondération d'activation des différents espaces.
On impose souvent une sparsité :
\[
|\supp(r(z,c))|\le k\ll n.
\]
\end{definition}

\begin{remark}
Cette sparsité formalise la gestion locale du calcul matriciel et de la mémoire : tous les espaces ne sont pas activés simultanément, seuls ceux jugés pertinents par le contexte participent à la computation.
\end{remark}

\subsection{Agents spécialisés}

\begin{definition}[Agent local]
Un agent local \(A_i\) est le triplet
\[
A_i=(\pi_i,\Pcal_i,\Theta_i),
\]
où \(\pi_i\) est la politique de décision locale, \(\Pcal_i\) l'opérateur de purification associé, et \(\Theta_i\) les paramètres internes de l'agent.
\end{definition}

Chaque agent est responsable d'un espace sémantique particulier, mais peut communiquer avec les autres agents via le graphe \(\Gcal\).

\begin{definition}[Mise à jour locale]
Si l'agent \(A_i\) est activé à l'instant \(t\), ses paramètres sont mis à jour par
\[
\Theta_i^{t+1}
=
\Theta_i^t
-
\eta_t \nabla_{\Theta_i}\mathcal L_i^t,
\]
tandis que les autres paramètres restent inchangés ou subissent une mise à jour nulle.
\end{definition}

\subsection{Traçabilité et déterminisme}

\begin{proposition}[Traçabilité]
Si la politique de routage \(r\), les opérateurs locaux \(\Pcal_i\), et l'opérateur de fusion \(\Gamma\) sont déterministes, alors la trajectoire \(\tau\) produite par le système PCH-ESM est elle-même déterministe pour une entrée et un contexte donnés.
\end{proposition}

\begin{remark}
Le déterminisme n'est donc pas automatique ; il doit être imposé. En revanche, l'architecture PCH-ESM le rend plus facile à contrôler qu'un grand espace sémantique monolithique.
\end{remark}

\subsection{Coherence locale et généralisation globale}

La spécialisation renforce la compétence locale mais fragilise la généralisation globale si aucune couche de coordination n'existe.

On distingue donc :

\begin{itemize}[leftmargin=2em]
    \item \textbf{généralisation locale} : bonne performance sur le domaine propre de l'espace spécialisé ;
    \item \textbf{généralisation globale} : bonne performance sur des tâches transversales, hors domaine initial, et sur des compositions multi-domaines.
\end{itemize}

L'architecture PCH-ESM seule ne garantit pas la généralisation globale. Celle-ci doit être reconstruite par la fonction \(G\).

% =================================================
\section{La fonction \(G\) : distillation de trajectoires et coordination inter-espaces}
% =================================================

On suppose l'existence :

\begin{itemize}[leftmargin=2em]
    \item d'un modèle généraliste \(M_G\) entraîné sur un corpus multi-disciplinaire ;
    \item d'une famille de modèles spécialisés \(\{M_i\}_{i=1}^n\), chacun entraîné sur un sous-domaine \(D_i\).
\end{itemize}

Le point clé n'est pas de distiller seulement des réponses, mais de distiller des trajectoires de coordination.

\subsection{Modèle généraliste}

Le modèle généraliste induit une politique de trajectoire
\[
\pi_G^{\mathrm{tr}}(\tau\mid x,c),
\]
où \(\tau\) est une trajectoire de pensée, de routage, de consultation d'experts et de fusion.

Le modèle généraliste apprend ainsi non seulement des connaissances, mais surtout des transitions inter-domaines.

\subsection{Modèles spécialisés}

Chaque modèle spécialisé \(M_i\) est entraîné sur un domaine \(D_i\).
Il doit être fort localement, c'est-à-dire performant sur son domaine, mais aussi apprendre à collaborer hors de son domaine.

On introduit un indicateur de domaine
\[
m_i(x)=
\begin{cases}
1 & \text{si } x\in D_i,\\
0 & \text{sinon.}
\end{cases}
\]

\subsection{Objectif de coordination}

L'objectif n'est pas que l'agent mathématique imite le comportement mathématique du généraliste sur son propre domaine, puisqu'il le maîtrise déjà localement.
L'objectif est qu'il apprenne \emph{quand et comment déléguer} ou coopérer lorsque le problème sort de son périmètre initial.

On sépare donc deux régimes :

\[
\text{si } m_i(x)=1,\quad \text{compétence locale forte};
\]
\[
\text{si } m_i(x)=0,\quad \text{coordination inter-spécialités}.
\]

\begin{definition}[Perte locale-coordination]
Pour le modèle \(M_i\), on définit
\[
\mathcal L_i
=
\mathbb E_{(x,y)\sim D_i}
\Bigl[
m_i(x)\,\mathcal L_i^{\mathrm{loc}}(x,y)
+
(1-m_i(x))\,\mathcal L_i^{\mathrm{coord}}(x)
\Bigr]
+
\lambda_i \mathcal L_i^{\mathrm{reg}}.
\]
\end{definition}

Ici :
\begin{itemize}[leftmargin=2em]
    \item \(\mathcal L_i^{\mathrm{loc}}\) mesure la compétence spécialisée ;
    \item \(\mathcal L_i^{\mathrm{coord}}\) mesure l'alignement avec les trajectoires du modèle généraliste hors domaine ;
    \item \(\mathcal L_i^{\mathrm{reg}}\) régularise le modèle ;
    \item \(\lambda_i>0\) pondère la régularisation.
\end{itemize}

\subsection{Définition de la fonction \(G\)}

\begin{definition}[Fonction de coordination \(G\)]
La fonction \(G\) est une application de distillation qui transforme les trajectoires générales en politiques locales coordonnées :
\[
G:\pi_G^{\mathrm{tr}}\longmapsto \{\pi_i^{\mathrm{coord}}\}_{i=1}^n.
\]
Plus précisément, \(G\) est obtenu par minimisation d'une divergence entre la trajectoire générale et la trajectoire produite par le système spécialisé :
\[
G^\star
=
\argmin_{G}
\;
\mathbb E_{(x,c)\sim \mathcal D}
\left[
\KL\!\left(
\pi_G^{\mathrm{tr}}(\cdot\mid x,c)
\;\middle\|\;
\pi_{G}^{\mathrm{spec}}(\cdot\mid x,c)
\right)
\right],
\]
sous contraintes de compétence locale sur chaque sous-domaine \(D_i\).
\end{definition}

\begin{remark}
La fonction \(G\) ne transfère pas la connaissance brute d'un modèle généraliste vers des spécialistes.
Elle transfère la \emph{forme des trajectoires} : quand consulter qui, dans quel ordre, avec quelle profondeur, et selon quelle hiérarchie de dépendances.
\end{remark}

\subsection{Fusion globale}

Une fois les politiques spécialisées alignées par \(G\), la sortie globale est obtenue par fusion :
\[
\hat y
=
D_\psi\!\Bigl(
\Gamma\bigl(h_1,\dots,h_n\bigr)
\Bigr),
\]
où \(h_i\) est la représentation locale issue de \(E_i\) et de sa purification propre.

\begin{proposition}[Principe de cohérence globale]
Si :
\begin{enumerate}[leftmargin=2em]
    \item chaque modèle spécialisé est compétent sur son domaine ;
    \item les politiques de coordination sont alignées avec la trajectoire générale via \(G\) ;
    \item la fusion \(\Gamma\) satisfait une condition de cohérence sur les espaces de recouvrement ;
\end{enumerate}
alors le système PCH-ESM coordonné par \(G\) peut approcher une généralisation globale comparable à celle d'un modèle généraliste, tout en conservant la spécialisation locale.
\end{proposition}

% =================================================
\section{Multi-PH : modèle du monde et états hiérarchiques}
% =================================================

On appelle Multi-PH la version où la purification hiérarchique ne sert plus seulement à construire une représentation utile pour une tâche, mais à modéliser l'évolution d'un \emph{état du monde} latent.

\begin{definition}[État du monde hiérarchique]
On note
\[
\Omega_t = \bigl(\omega_t^{(1)},\dots,\omega_t^{(L)}\bigr)
\]
l'état du monde à l'instant \(t\), organisé en niveaux hiérarchiques.
Chaque niveau \(\omega_t^{(\ell)}\) représente une abstraction différente du monde.
\end{definition}

\begin{definition}[Dynamique PCH du monde]
La dynamique de monde est donnée par
\[
\Omega_{t+1}
=
F\!\bigl(\Omega_t,a_t,o_t\bigr),
\]
où \(a_t\) est une action, \(o_t\) une observation, et \(F\) est une application construite elle-même à partir d'une purification contextuelle hiérarchique.
\end{definition}

\begin{remark}
Dans cette perspective, prédire le prochain token n'est qu'un cas particulier d'une observation issue de l'état du monde.
La vraie variable à prédire est l'état hiérarchique \(\Omega_{t+1}\), et le token est un readout.
\end{remark}

\begin{definition}[Readout token]
Le prochain token est obtenu par une tête de lecture
\[
p_\psi(x_{t+1}\mid \Omega_{t+1})
=
D_\psi(\Omega_{t+1}).
\]
\end{definition}

% =================================================
\section{Architecture unifiée}
% =================================================

L'architecture complète peut être schématisée par
\[
x_{1:T}
\longrightarrow
E_G(x_{1:T})
\longrightarrow
Z_G
\longrightarrow
\{\E_i\}_{i=1}^n
\longrightarrow
\{\tau_i\}_{i=1}^n
\longrightarrow
\Gamma
\longrightarrow
\hat y,
\]
ou, dans la variante Multi-PH,
\[
x_{1:T}
\longrightarrow
\Omega_{1:T}
\longrightarrow
p_\psi(x_{t+1}\mid \Omega_{t+1}).
\]

\subsection{Fonction de perte globale}

On peut écrire une perte globale de la forme
\[
\mathcal L_{\mathrm{tot}}
=
\lambda_{\mathrm{task}}\mathcal L_{\mathrm{task}}
+
\lambda_{\mathrm{pch}}\mathcal L_{\mathrm{pch}}
+
\lambda_{\mathrm{coord}}\mathcal L_{\mathrm{coord}}
+
\lambda_{\mathrm{multiES}}\mathcal L_{\mathrm{multiES}}
+
\lambda_{\mathrm{world}}\mathcal L_{\mathrm{world}}
+
\lambda_{\mathrm{reg}}\Omega(\Theta,\Psi),
\]
où :
\begin{itemize}[leftmargin=2em]
    \item \(\mathcal L_{\mathrm{task}}\) est la perte de tâche ;
    \item \(\mathcal L_{\mathrm{pch}}\) force la stabilité contextuelle ;
    \item \(\mathcal L_{\mathrm{coord}}\) aligne les trajectoires entre espaces ;
    \item \(\mathcal L_{\mathrm{multiES}}\) pénalise l'incohérence entre espaces ;
    \item \(\mathcal L_{\mathrm{world}}\) entraîne le modèle du monde ;
    \item \(\Omega\) est une régularisation standard.
\end{itemize}

% =================================================
\section{Interprétation théorique}
% =================================================

La théorie PCH-ESM affirme les points suivants :

\begin{enumerate}[leftmargin=2em]
    \item La bonne représentation n'est pas seulement celle qui prédit, mais celle qui reste stable sous purification récursive dans un contexte donné.
    \item L'attention standard peut être comprise comme un cas plat de ce principe.
    \item Les poids eux-mêmes peuvent être traités comme mémoire purifiable.
    \item Un espace sémantique monolithique peut être remplacé par une hiérarchie d'espaces spécialisés.
    \item La généralisation globale ne disparaît pas nécessairement avec la spécialisation locale ; elle peut être reconstruite par une fonction de coordination \(G\) apprenant les trajectoires d'un modèle généraliste.
    \item Les trajectoires internes de calcul deviennent explicites, donc traçables et potentiellement déterministes.
    \item La version monde du formalisme (Multi-PH) remplace la prédiction de tokens comme objectif fondamental par la prédiction d'états hiérarchiques latents.
\end{enumerate}

\begin{remark}
Cette théorie ne suppose pas que la spécialisation soit gratuite : elle exige une couche de coordination pour préserver la généralisation transversale.
C'est précisément le rôle de \(G\).
\end{remark}

\section{Conclusion}

PCH-ESM fournit une formalisation unifiée de :
\begin{itemize}[leftmargin=2em]
    \item la purification contextuelle hiérarchique ;
    \item la version adaptative des paramètres ;
    \item la hiérarchisation de l'espace sémantique ;
    \item la décentralisation de l'IA en espaces spécialisés ;
    \item la coordination inter-espaces par apprentissage de trajectoires ;
    \item le modèle du monde hiérarchique.
\end{itemize}

La proposition centrale est que la structure cognitive utile n'est pas uniquement un grand espace latent global, mais un \emph{écosystème hiérarchique d'espaces sémantiques} dont la stabilité, la spécialisation et la coordination sont explicitement représentées.

\end{document}